% =================================================================
% Sprint 4 - section 3 feature description (kvy)
% AI content moderation for forum posts and AI face scan / identity
% extraction with duplicate detection.
% =================================================================

% =================================================================
% 3.x AI Content Moderation
% Feature description written against the real implementation:
%   apps/api/src/shared/services/ai-moderation.service.ts
%   apps/api/src/routes/forums/forums-post/forums.service.ts
%   apps/web/app/forum/[categorySlug]/components/create-post-dialog.tsx
%   apps/web/app/admin/moderation/* and apps/web/app/admin/trash/*
% Procedures are described from the user interface point of view.
% =================================================================

\Needspace{12\baselineskip}
\subsection{AI Content Moderation}

\Needspace{12\baselineskip}
\subsubsection{Moderate Content by AI}

\functionfield{Function Trigger} A logged-in Client or Freelancer opens a forum category, clicks \textbf{New Post}, fills in the title and content of the post and presses the \textbf{Post} button in the create-post dialog.

\functionfield{Function Description} This function automatically scores the content of a forum post with an AI moderation service when the post is submitted. The title and content are sent to \texttt{AIModerationService} (via \texttt{AI\_MODERATION\_URL}/moderate), which returns a score, detected categories, and a moderation verdict (\texttt{BLOCK} mapped to \texttt{REJECTED}, \texttt{REVIEW} to \texttt{PENDING}, and \texttt{ALLOW} to \texttt{APPROVED}). Regardless of the verdict, the post is persisted in PostgreSQL with its moderation status: \texttt{REJECTED} posts stay hidden from public listings and appear in the administrator trash (with HTTP 201 returned and an error toast shown in the create dialog); \texttt{PENDING} posts wait in the moderation queue; and \texttt{APPROVED} posts become publicly visible immediately. When the external AI moderation service is unavailable, the backend defaults to \texttt{APPROVED} so posting is not blocked.

\screenfield[0.92]{image/screenshot/sprint4/kvy/CreatePostDialog.png}{Create post dialog with title and rich-text content}
\screenfield[0.92]{image/screenshot/sprint4/kvy/AdminModeration.png}{Administrator moderation queue for posts that need manual review}
\screenfield[0.92]{image/screenshot/sprint4/kvy/AdminTrash.png}{Trash list with posts rejected during moderation}

\functionfield{Function Details}
\begin{itemize}[leftmargin=1.6em,topsep=4pt,itemsep=3pt,parsep=0pt,label=$\bullet$]
    \item \detailhead{Precondition}
        \begin{itemize}[leftmargin=2.2em,topsep=2pt,itemsep=3pt,parsep=0pt,label=$\circ$]
            \item The user is logged in with a Client or Freelancer account.
            \item The user is connected to the internet.
            \item The forum category exists and the user is allowed to create a post in it.
            \item The title and the content of the post are not empty.
        \end{itemize}
    \item \detailhead{Process}
        \begin{itemize}[leftmargin=2.2em,topsep=2pt,itemsep=3pt,parsep=0pt,label=$\circ$]
            \item \textbf{Open the Create Post Dialog}: The user opens a forum category and clicks \textbf{New Post}, then types a title and the content in the dialog.
            \item \textbf{Submit the Post}: The user presses the \textbf{Post} button. The frontend sends the request to create the post to the backend.
            \item \textbf{Run the AI Check}: The backend calls the AI moderation service with the combined title and content of the post and the context \texttt{forum-post} (falling back to \texttt{APPROVED} if unreachable).
            \item \textbf{Receive the Result}: The service returns a score, detected categories, a decision (\texttt{BLOCK}/\texttt{REVIEW}/\texttt{ALLOW}), and a message.
            \item \textbf{Persist Verdict and Apply Visibility}: The backend maps the verdict and creates the \texttt{ForumPost} record in PostgreSQL, returning HTTP 201.
                \begin{itemize}[leftmargin=2.2em,topsep=2pt,itemsep=3pt,parsep=0pt,label=$\blacksquare$]
                    \item \textbf{Rejected (\texttt{BLOCK} $\rightarrow$ \texttt{REJECTED})}: The post is saved with status \texttt{REJECTED} (hidden from public listings and visible in the administrator trash). The dialog stays open and shows an error toast explaining that the content was rejected.
                    \item \textbf{Pending (\texttt{REVIEW} $\rightarrow$ \texttt{PENDING})}: The post is saved with status \texttt{PENDING}. The dialog closes, the form is cleared and an information message tells the user that the post is waiting for administrator review. The post is not shown in public listings.
                    \item \textbf{Approved (\texttt{ALLOW} $\rightarrow$ \texttt{APPROVED})}: The post is saved with status \texttt{APPROVED}. The dialog closes, the form is cleared and a success message is shown. The user is taken to the detail page of the new post, which is immediately visible to everyone.
                \end{itemize}
            \item \textbf{Review by an Administrator}: When a post is \texttt{PENDING}, it appears in the administrator moderation queue, where it can be approved or rejected (see the sprint 3 moderation features).
            \item \textbf{View Rejected Posts in Trash}: Both AI-\texttt{REJECTED} posts and administrator-rejected or soft-deleted posts appear in the administrator trash, where they can be restored.
        \end{itemize}
    \item \detailhead{Alternative Flow}
        \begin{itemize}[leftmargin=2.2em,topsep=2pt,itemsep=3pt,parsep=0pt,label=$\circ$]
            \item \textbf{Empty Title or Content}: The \textbf{Post} button stays disabled while the title or the content is empty, so no request is sent.
            \item \textbf{Cancel the Dialog}: When the user closes the dialog or presses \textbf{Cancel}, no moderation request is made and nothing is saved.
            \item \textbf{AI Service Unavailable}: When the AI service cannot be reached, \texttt{AIModerationService} defaults to \texttt{APPROVED} so users are not blocked, and the post is saved with status \texttt{APPROVED}.
            \item \textbf{Lost Connection}: When the connection to the server is lost while submitting, the post is not created and the dialog shows an error message.
        \end{itemize}
    \item \detailhead{Postconditions}
        \begin{itemize}[leftmargin=2.2em,topsep=2pt,itemsep=3pt,parsep=0pt,label=$\circ$]
            \item An AI-\texttt{REJECTED} post is saved with status \texttt{REJECTED}, hidden from public listings, and visible in the administrator trash.
            \item A pending post is stored with status \texttt{PENDING}, is absent from public listings and appears in the moderation queue.
            \item An approved post is stored with status \texttt{APPROVED} and is visible in the forum immediately.
            \item A post rejected by an administrator is hidden from the public and can be found in the trash for restoration.
        \end{itemize}
    \item \detailhead{Error Handling}
        \begin{itemize}[leftmargin=2.2em,topsep=2pt,itemsep=3pt,parsep=0pt,label=$\circ$]
            \item \textbf{Not Authenticated}: When a guest tries to create a post, the user is sent to the login page.
            \item \textbf{Content Rejected}: When the AI verdict is \texttt{BLOCK}, the post is persisted as \texttt{REJECTED} and the frontend keeps the dialog open showing the moderation error message.
            \item \textbf{Validation Error}: When the title or the content does not pass validation, the backend returns HTTP 400 and the form shows the corresponding error.
            \item \textbf{Lost Connection}: When the connection to the server is lost, the post is not created and the dialog shows an error message.
        \end{itemize}
    \item \detailhead{Security Measures}
        \begin{itemize}[leftmargin=2.2em,topsep=2pt,itemsep=3pt,parsep=0pt,label=$\circ$]
            \item Only authenticated users may create forum posts, and the check is made again on the server.
            \item Moderation verdicts are resolved on the server and cannot be bypassed from the browser.
            \item Pending posts are never returned to ordinary users, because public listings only contain posts with status \texttt{APPROVED}.
            \item The moderation queue and the trash are only available to administrators.
        \end{itemize}
\end{itemize}

% =================================================================
% 3.x AI Face Scan & Duplicate Detection
% Feature description written against the real implementation:
%   apps/api/src/routes/account-profile/account-profile.service.ts
%   apps/web/app/(freelancer)/freelancer/profile/identity-verification-form.tsx
%   apps/web/app/admin/identity-verifications/*
% Procedures are described from the user interface point of view.
% =================================================================

\Needspace{12\baselineskip}
\subsection{AI Face Scan \& Duplicate Detection}

\Needspace{12\baselineskip}
\subsubsection{Scan Face and Detect Duplicate Identity}

\functionfield{Function Trigger} A logged-in Freelancer or Client opens \textbf{Profile} $\rightarrow$ \textbf{Identity} tab, uploads the required identity documents and selfie, completes the liveness check and submits the verification request.

\functionfield{Function Description} This function uses an AI service (\texttt{id-verification} with RapidOCR, UltraFace/ArcFace ONNX, and MediaPipe) to verify the face from the uploaded images, extract OCR text and document number (\texttt{document\_number}), check liveness, and store the verification document in PostgreSQL. When an extracted document number matches an already-approved record of another user, the submission is recorded with status \texttt{REJECTED} and linked via \texttt{duplicateOfId}; otherwise, an AI \texttt{VERIFIED} outcome is saved as \texttt{PENDING} for administrator review, and an AI \texttt{REJECTED} or \texttt{RETRY} outcome is saved as \texttt{REJECTED}. Users can follow the request status from the Identity tab, and administrators review pending submissions from the identity verification console.

\screenfield[0.92]{image/screenshot/sprint4/kvy/IdentityVerificationTab.png}{Freelancer profile with the identity verification tab}
\screenfield[0.92]{image/screenshot/sprint4/kvy/IdentityLivenessCheck.png}{Liveness challenge while scanning the face}
\screenfield[0.92]{image/screenshot/sprint4/kvy/AdminIdentityVerifications.png}{Administrator screen for reviewing identity verification requests}

\functionfield{Function Details}
\begin{itemize}[leftmargin=1.6em,topsep=4pt,itemsep=3pt,parsep=0pt,label=$\bullet$]
    \item \detailhead{Precondition}
        \begin{itemize}[leftmargin=2.2em,topsep=2pt,itemsep=3pt,parsep=0pt,label=$\circ$]
            \item The user is logged in with a Freelancer or Client account.
            \item The user is connected to the internet.
            \item The front image of the identity document is provided, together with the selfie and the liveness frames required by the form.
            \item The user does not already hold an \texttt{APPROVED} identity verification record.
        \end{itemize}
    \item \detailhead{Process}
        \begin{itemize}[leftmargin=2.2em,topsep=2pt,itemsep=3pt,parsep=0pt,label=$\circ$]
            \item \textbf{Open the Identity Tab}: The user opens their profile settings and selects the \textbf{Identity} tab, which shows the current verification status and the upload form.
            \item \textbf{Upload the Documents}: The user selects the document type and uploads the front image of the document, and optionally the back image.
            \item \textbf{Capture the Selfie and Liveness}: The user takes a selfie and completes the liveness challenge, which captures the frames used for the face scan.
            \item \textbf{Submit the Request}: The user presses the submit button. The frontend sends the document type, the images and the liveness frames to \texttt{POST /api/identity-verifications/ai-check}.
            \item \textbf{Validate the Files}: The backend checks that the required images are present, match JPEG/PNG/WebP formats, and do not exceed 10 MB.
            \item \textbf{Scan the Face and Extract Document Number}: The backend calls the AI identity service, which runs RapidOCR to extract raw text and \texttt{document\_number}, compares the card photo and selfie with UltraFace/ArcFace, and verifies liveness poses with MediaPipe.
            \item \textbf{Check Duplicate Identity Number}: Using any extracted \texttt{document\_number}, the backend queries existing \texttt{APPROVED} verification records of other users via \texttt{findApprovedIdentityDocumentByIdentityNumber}.
            \item \textbf{Store the Result}: Uploaded images are persisted and an \texttt{IdVerificationDocument} row is created with status \texttt{PENDING} (when AI returns \texttt{VERIFIED} and no duplicate is found) or \texttt{REJECTED} (when AI returns \texttt{REJECTED}/\texttt{RETRY} or a duplicate \texttt{identityNumber} matches another account, recording \texttt{duplicateOfId}).
            \item \textbf{Review by an Administrator}: Submissions saved as \texttt{PENDING} appear in the administrator identity verification screen, where an administrator reviews the uploaded images and approves or rejects the request.
            \item \textbf{Show the Status}: The user sees the current status of the request in the Identity tab (\texttt{PENDING}, \texttt{APPROVED}, or \texttt{REJECTED} with notes).
        \end{itemize}
    \item \detailhead{Alternative Flow}
        \begin{itemize}[leftmargin=2.2em,topsep=2pt,itemsep=3pt,parsep=0pt,label=$\circ$]
            \item \textbf{Missing Required File}: When the front image, the selfie or the liveness frames are missing, the submission is rejected and the form asks the user to complete them.
            \item \textbf{Unsupported File}: When an uploaded file is not an accepted image type or is too large (above 10 MB), the form or backend rejects it and no AI scan runs.
            \item \textbf{AI Rejects or Requests Retry}: When the AI returns \texttt{REJECTED} or \texttt{RETRY}, the backend stores the document with status \texttt{REJECTED} and returns the rejection reason so the user can resubmit clearer photos.
            \item \textbf{Duplicate Detected}: When the extracted document number matches another account's approved record, the submission is recorded with status \texttt{REJECTED} and \texttt{duplicateOfId} referencing the matched record.
            \item \textbf{AI Service Unavailable}: When the AI service cannot be reached (HTTP 503/504), no verification document is created in PostgreSQL and the user is asked to try again later.
        \end{itemize}
    \item \detailhead{Postconditions}
        \begin{itemize}[leftmargin=2.2em,topsep=2pt,itemsep=3pt,parsep=0pt,label=$\circ$]
            \item On AI \texttt{VERIFIED} without duplicates, a verification record is stored with status \texttt{PENDING} alongside \texttt{aiSummary} and \texttt{aiSelfieSummary}.
            \item On AI \texttt{REJECTED}/\texttt{RETRY} or duplicate detection, a verification record is stored with status \texttt{REJECTED} (and \texttt{duplicateOfId} when applicable).
            \item The administrator can review, approve or reject pending requests from the identity verification screen.
            \item The user can see the current status of the request from the Identity tab.
        \end{itemize}
    \item \detailhead{Error Handling}
        \begin{itemize}[leftmargin=2.2em,topsep=2pt,itemsep=3pt,parsep=0pt,label=$\circ$]
            \item \textbf{Not Authenticated}: When a guest tries to open the identity tab, the user is sent to the login page.
            \item \textbf{Missing or Invalid File}: When a required file is missing or invalid, the backend returns an error and the form shows the corresponding message.
            \item \textbf{AI Service Error or Timeout}: When the AI service is unreachable or times out, the backend throws an error without creating a verification record.
            \item \textbf{Lost Connection}: When the connection to the server is lost during submission, the request is not saved and the user is asked to try again.
        \end{itemize}
    \item \detailhead{Security Measures}
        \begin{itemize}[leftmargin=2.2em,topsep=2pt,itemsep=3pt,parsep=0pt,label=$\circ$]
            \item Only the owner of the profile may submit identity documents, and the check is made again on the server.
            \item Uploaded files are validated for type and size before they are processed.
            \item Extracted identity data and file references are stored so that the request can be audited.
            \item Duplicate detection is performed on the server, so repeated registrations cannot be hidden from the browser.
            \item The identity verification screens are only available to administrators.
        \end{itemize}
\end{itemize}
